\section{Data}\label{sec:data}

\subsection{Policy documents and PCSI}

The policy data come from the corpus assembled for the companion paper \citep{SharmaPCSI2026}: 519 IP and patent policy
documents from 150 U.S.\ research universities, dated 1944--2025, forming 481 institution-year document records and 87,160
sentences. PCSI is the mean sentence-level stance score of a document, where each sentence is scored from restrictive to
supportive toward inventors; its construction and validation are described in the companion paper. For the annual panel,
each observed document is treated as the policy in force until the next observed document.

\subsection{Technology-transfer outcomes}

Outcomes come from the AUTM Licensing Activity Survey for 1991--2023. The raw data contain 253 reporting identifiers,
which we harmonize to 149 institutions. Harmonization matters because 101 institutions report under more than one
identifier over the period; treating identifiers as institutions splits institution fixed effects and breaks lagged
variables whenever an identifier changes. All analyses are keyed to the harmonized institution.

The principal outcomes are licenses and options executed, invention disclosures received, new U.S.\ patent applications
filed, and U.S.\ patents issued, each entered as $\log(1+y)$, together with the filing margin
\begin{equation}
F_{it}=\log(1+\text{Applications}_{it})-\log(1+\text{Disclosures}_{it}),
\end{equation}
an institution-year analogue of the filing ratio \citep{ChoudhryPonzio2020}. Because the survey does not link individual
disclosures to applications, $F_{it}$ is not an invention-level conversion probability.

Two outcome series required auditing. In the FY2023 survey the total licenses-and-options field no longer equals licenses
issued plus options issued, as it does in 99 percent of 2005--2022 records; its 2023 median is 6 against 29 for the
component sum. We therefore use the component sum for 2023 so that the series keeps one definition (Appendix
Table~\ref{tab:licfield}). Startups formed are reported only from 1994; in the survey's definition they are companies formed
specifically to license the institution's technology, and counts are self-reported. The series is zero for 39 percent of
reporting institutions in 1994--1999 but only 8 percent in 2015--2022, rising to 17 percent in 2023, and between 2 and 17
institutions enter or leave the set of reporters each year. We found no evidence of miscoded zeros, but the short coverage,
the change in the zero share, and the known variation in how institutions apply the definition make startups less
comparable over time than the other outcomes. We therefore report startups only as supplementary evidence
(Appendix~\ref{app:startups}). Licensing income is not used because its reported units are inconsistent across survey years.

Table~\ref{tab:descriptives} summarizes the linked panel. In the annual panel, 89 percent of PCSI values are carried
forward from an earlier document, and 30 percent of PCSI variance is within institutions.

\begin{table}[htbp]\centering
\caption{Linked AUTM--PCSI panel, 1991--2023}\label{tab:descriptives}
\begin{threeparttable}\small\begin{tabular}{lrrr}\toprule
Variable & $N$ & Mean & SD \\\midrule
PCSI (policy in force) & 2,564 & 0.440 & 0.063 \\
Invention disclosures & 3,507 & 122.4 & 157.7 \\
New patent applications & 3,507 & 70.6 & 113.7 \\
U.S. patents issued & 3,351 & 33.0 & 50.3 \\
Licenses and options executed$^{a}$ & 3,489 & 40.0 & 69.2 \\
Research expenditure (\$ millions) & 3,507 & 333.0 & 478.1 \\
Licensing FTEs & 3,401 & 5.5 & 7.3 \\
Startups formed$^{b}$ & 3,245 & 4.6 & 7.1 \\
\midrule
Institution-years & \multicolumn{3}{r}{3,507} \\
Harmonized institutions (AUTM reporting identifiers) & \multicolumn{3}{r}{149 (253)} \\
Share of annual PCSI values carried forward & \multicolumn{3}{r}{0.893} \\
Within-institution share of PCSI variance & \multicolumn{3}{r}{0.30} \\
\bottomrule\end{tabular}
\begin{tablenotes}[flushleft]\footnotesize
\item \textit{Notes:} Institution-year observations. PCSI is the policy-in-force score from the companion measurement paper, carried forward from the most recent observed document. $^{a}$For FY2023 the survey's total field no longer equals licenses plus options issued (as it does in 99\% of 2005--2022 records), so 2023 values use licenses plus options issued. $^{b}$Supplementary outcome, reported from 1994 (Appendix~\ref{app:startups}). Licensing income is not used because its reported units are inconsistent across survey years.
\end{tablenotes}\end{threeparttable}\end{table}

\subsection{Revisions and the revision sample}

A revision is an observed policy document whose PCSI differs from the institution's previous observed document by more
than 0.03, approximately one-half of the cross-document standard deviation of the index; Appendix
Table~\ref{tab:thresholduniverse} shows how the number of detected revisions changes at thresholds from 0.02 to 0.05.
Revisions are identified from the full sequence of observed policy documents rather than from the AUTM-linked annual
panel, because documents issued in years without AUTM observations would otherwise be skipped. The corpus contains the
documents that were collected, not every version ever adopted, so documentary review is essential.

At the 0.03 threshold there are 127 revisions at 78 institutions linked to AUTM outcomes. We separate sample construction
into a mechanical screen and a documentary screen. The mechanical screen requires that the institution have no other
threshold revision in the event window and that core AUTM outcomes provide usable observations around the event.
Sixty-two revisions pass these restrictions. Twenty-five had already been hand-reviewed in the September-25 manuscript.
We subsequently reviewed the predecessor and successor documents for all 37 remaining mechanically eligible pairs, so
every revision capable of entering the preferred event-study sample is now subject to the same documentary standard.

Each pair is classified on two dimensions. \emph{Comparability} asks whether the two documents are the same governing IP
or patent policy (C), partially comparable (P: for example, an excerpt compared with a full policy or a campus policy
compared with a governing-board policy), or not comparable (N). \emph{Timing} asks whether the revision can be dated.
Tier A requires a stated adoption or revision date with no documented intervening version, or documents no more than two
years apart; tier B has a stated adoption or effective date but no revision history; tier C includes pairs for which
intervening versions are documented or no usable revision date can be established. The event year is the documented
adoption or revision year rather than the filename year whenever those differ.

Of the 37 newly reviewed pairs, nine satisfy C/P comparability and A/B timing; 28 fail at least one documentary
requirement. Appending those nine to the original hand-reviewed events yields a preferred sample of 34 revisions at 32
institutions: eight upward and 26 downward. All 34 therefore satisfy the same stated documentary rules. The main
specification allows an event when an outcome is observed in the event year even if the AUTM series ends immediately
afterward. Because this matters for late events and sparse reporters, we also define a stricter sample requiring at least
two treated-institution core-outcome years in event times 0 through $+5$, which guarantees at least one observed calendar
year after the revision. That rule yields 29 revisions, seven upward and 22 downward. We report it as a uniform sensitivity
rather than selectively removing late events.

\begin{table}[htbp]\centering
\caption{Construction of the revision sample}\label{tab:revsample}
\begin{threeparttable}\small
\begin{tabular}{p{11.3cm}r}\toprule
Step & Count \\\midrule
Observed policy-document records (institution-years) in the companion corpus & 481 \\
Revisions with $|\Delta\mathrm{PCSI}|>0.03$ at AUTM-linked institutions & 127 \\
Pass mechanical non-overlap and outcome-support screen & 62 \\
\quad previously hand-reviewed in September-25 specification & 25 \\
\quad additional mechanically eligible pairs reviewed in the second pass & 37 \\
Newly usable C/P pairs with A/B timing from the second pass & 9 \\
Preferred fully documentary-reviewed sample & 34 \\
\quad upward / downward & 8 / 26 \\
\quad distinct revising institutions & 32 \\
Stricter post-support sensitivity sample & 29 \\
\quad upward / downward & 7 / 22 \\
\bottomrule
\end{tabular}
\begin{tablenotes}[flushleft]\footnotesize
\item \textit{Notes:} Mechanical eligibility is determined without using estimated outcome effects. Documentary usability
requires C or P comparability and timing tier A or B. The 65 revisions failing the mechanical screen cannot enter the
preferred stacked event study regardless of documentary classification. The stricter 29-event sample requires at least
two treated-institution core-outcome panel years in event times 0--5.
\end{tablenotes}
\end{threeparttable}
\end{table}
